%%!TeX program = xelatex
%% 부득이하게 pdflatex을 사용해야 할 경우 위의 magic comment를 제거하십시오.

% Initiated by 정민석(2014년도 경기과학고 수학과전문교원)
% Continously being modified by 경기과학고 TeX 사용자협회
% Website : http://gshslatexintro.github.io 
% 박기현(선생님)과 한수민(37기) 학생이 본문 나누어진 버전으로 개선함.
% Website : http://gshslatexintro.github.io 


\documentclass{gshs_thesis}
\RequirePackage{xetexko}
\setmainfont{Times New Roman}
\setmainhangulfont{Batang}
\nonfrenchspacing
\graphicspath{{images/}}
% 이곳에 필요한 별도의 패키지들을 적어넣으시오.
%\usepackage{...}
\usepackage{verbatim} % for commment, verbatim environment
\usepackage{spverbatim} % automatic linebreak verbatim environment
\usepackage{listings}
\lstset{
	basicstyle=\small\ttfamily,
	columns=flexible,
	breaklines=true
}

% -----------------------------------------------------------------------
%                   이 부분은 수정하지 마시오.
% -----------------------------------------------------------------------
\titleheader{졸업논문청구논문}
\school{과학영재학교 경기과학고등학교}
\approval{위 논문은 과학영재학교 경기과학고등학교 졸업논문으로\\
졸업논문심사위원회에서 심사 통과하였음.}
\chairperson{심사위원장}
\examiner{심사위원}
\apprvsign{(인)}
\korabstract{초 록}
\koracknowledgement{감사의 글}
\korresearches{연 구 활 동}

%: ----------------------------------------------------------------------
%:                  논문 제목과 저자 이름을 입력하시오
% ----------------------------------------------------------------------
\title{한글 제목} %한글 제목
\engtitle{English Title} %영문 제목
\korname{홍 길 동} %저자 이름을 한글로 입력하시오 (글자 사이 띄어쓰기)
\engname{Hong, Gil-Dong} %저자 이름을 영어로 입력하시오 (family name, personal name)
\chnname{洪 吉 東} %저자 이름을 한자로 입력하시오 (글자 사이 띄어쓰기)
\studid{14201} %학번을 입력하시오

%------------------------------------------------------------------------
%                  심사위원과 논문 승인 날짜를 입력하시오
%------------------------------------------------------------------------
\advisor{Mok, Chinook}  %지도교사 영문 이름 (family name, personal name)
\judgeone{박 승 원} %심사위원장
\judgetwo{김 대 감}   %심사위원1
\judgethree{목 진 욱} %심사위원2(지도교사)
\degreeyear{2017}   %졸업 년도
\degreedate{2016}{11}{13} %논문 승인 날짜 양식

%------------------------------------------------------------------------
%                  논문제출 전 체크리스트를 확인하시오
%------------------------------------------------------------------------
\checklisttitle{[논문제출 전 체크리스트]} %수정하지 마시오
\checklistI{1. 이 논문은 내가 직접 연구하고 작성한 것이다.} %수정하지 마시오
% 이 항목이 사실이라면 다음 줄 앞에 "%"기호 삽입, 다다음 줄 앞의 "%"기호 제거하시오
\checklistmarkI{$\square$}
%\checklistmarkI{$\text{\rlap{$\checkmark$}}\square$}
\checklistII{2. 인용한 모든 자료(책, 논문, 인터넷자료 등)의 인용표시를 바르게 하였다.} %수정하지 마시오
% 이 항목이 사실이라면 다음 줄 앞에 "%"기호 삽입, 다다음 줄 앞의 "%"기호 제거하시오
\checklistmarkII{$\square$}
%\checklistmarkII{$\text{\rlap{$\checkmark$}}\square$}
\checklistIII{3. 인용한 자료의 표현이나 내용을 왜곡하지 않았다.} %수정하지마시오
% 이 항목이 사실이라면 다음 줄 앞에 "%"기호 삽입, 다다음 줄 앞의 "%"기호 제거하시오
\checklistmarkIII{$\square$}
%\checklistmarkIII{$\text{\rlap{$\checkmark$}}\square$}
\checklistIV{4. 정확한 출처제시 없이 다른 사람의 글이나 아이디어를 가져오지 않았다.} %수정하지 마시오
% 이 항목이 사실이라면 다음 줄 앞에 "%"기호 삽입, 다다음 줄 앞의 "%"기호 제거하시오
\checklistmarkIV{$\square$}
%\checklistmarkIV{$\text{\rlap{$\checkmark$}}\square$}
\checklistV{5. 논문 작성 중 도표나 데이터를 조작(위조 혹은 변조)하지 않았다.} %수정하지 마시오
% 이 항목이 사실이라면 다음 줄 앞에 "%"기호 삽입, 다다음 줄 앞의 "%"기호 제거하시오
\checklistmarkV{$\square$}
%\checklistmarkV{$\text{\rlap{$\checkmark$}}\square$}
\checklistVI{6. 다른 친구와 같은 내용의 논문을 제출하지 않았다.} %수정하지 마시오
% 이 항목이 사실이라면 다음 줄 앞에 "%"기호 삽입, 다다음 줄 앞의 "%"기호 제거하시오
\checklistmarkVI{$\square$}
%\checklistmarkVI{$\text{\rlap{$\checkmark$}}\square$} % usepackage 등의 명령어는 여기에.
\usepackage{cite}
\usepackage{textcomp}
\usepackage{tocloft}
\setlength{\cftbeforesecskip}{0pt}
\setlength{\cftbeforesubsecskip}{0pt}
\setlength{\cftbeforesubsubsecskip}{0pt}

% 본문 시작
\begin{document}

	%표지만들기
	%makecover 함수와 관련하여 "Underfull \hbox (badness 10000) in paragraph" 오류는 무시하십시오. (TeXstudio ver 2.9.4 오류 기준)
	%\makecover
	

	%\baselineskip=2.2em         % line spacing in the paragraph
	\renewcommand\baselinestretch{1.2} % line spacing in the paragraph
\baselineskip=22pt plus1pt         % line spacing in the paragraph

\maketitle  % command to print the title page with above variables
\setcounter{page}{1}
%---------------------------------------------------------------------
%                  영문 초록을 입력하시오
%---------------------------------------------------------------------
\begin{abstracts}     %this creates the heading for the abstract page
\noindent{
Put your abstract here. It is completely consistent with 한글초록.
}
\end{abstracts}

%---------------------------------------------------------------------
%                  국문 초록을 입력하시오
%---------------------------------------------------------------------
\begin{abstractskor}        %this creates the heading for the abstract page
\noindent{
초록(요약문)은 가장 마지막에 작성한다. 연구한 내용, 즉 본론부터 요약한다. 서론 요약은 하지 않는다. 대개 첫 문장은 연구 주제 (+방법을 핵심적으로 나타낼 수 있는 문구: 실험적으로, 이론적으로, 시뮬레이션을 통해)를 쓴다. 다음으로 연구 방법을 요약한다. 선행 연구들과 구별되는 특징을 중심으로 쓴다. 뚜렷한 특징이 없다면 연구방법은 안써도 상관없다. 다음으로 연구 결과를 쓴다. 연구 결과는 추론을 담지 않고, 객관적으로 서술한다. 마지막으로 결론을 쓴다. 이 연구를 통해 주장하고자 하는 바를 간략히 쓴다. 요약문 전체에서 연구 결과와 결론이 차지하는 비율이 절반이 넘도록 한다. 읽는 이가 요약문으로부터 얻으려는 정보는 연구 결과와 결론이기 때문이다. 연구 결과만 레포트하는 논문인 경우, 결론을 쓰지 않는 경우도 있다.
}
\end{abstractskor}


%----------------------------------------------
%   Table of Contents (자동 작성됨)
%----------------------------------------------
\setcounter{secnumdepth}{3} % organisational level that receives a numbers
\setcounter{tocdepth}{3}    % print table of contents for level 3
\tableofcontents


%----------------------------------------------
%     List of Figures/Tables (자동 작성됨)
%----------------------------------------------
\cleardoublepage
\clearpage
\listoffigures	% 그림 목록과 캡션을 출력한다. 만약 논문에 그림이 없다면 이 줄의 맨 앞에 %기호를 넣어서 코멘트 처리한다.

\cleardoublepage
\clearpage
\listoftables  % 표 목록과 캡션을 출력한다. 만약 논문에 표가 없다면 이 줄의 맨 앞에 %기호를 넣어서 코멘트 처리한다.

\cleardoublepage
\clearpage
\renewcommand{\thepage}{\arabic{page}}
\setcounter{page}{1} % Abstract
	
	%%%%%%%%%%%%%%%%%%%%%%%%%%%%%%%%%%%%%%%%%%%%%%%%%%%%%%%%%%%
	%%%% Main Document %%%%%%%%%%%%%%%%%%%%%%%%%%%%%%%%%%%%%%%%
	%%%%%%%%%%%%%%%%%%%%%%%%%%%%%%%%%%%%%%%%%%%%%%%%%%%%%%%%%%%
	\cleardoublepage
	\clearpage
	\renewcommand{\thepage}{\arabic{page}}
	\setcounter{page}{1}
	
	%각 장을 아래와 같이 sub 폴더 안에 만들어서 넣으면 편리하다.
	\section{Introduction}

서론은 연구를 진행하게 된 배경을 기술하는 곳으로 보통 다음과 같은 순서로 쓰는 편이다.
\begin{itemize}
\item{연구 주제의 전반적 관심을 조명.}
\item{연구 분야의 스페셜 이슈를 조명.}
\item{해당 이슈를 해결하기 위한 다양한 선행 연구들을 서술.}
\item{선행 연구들의 한계점을 기술.}
\item{한계를 극복하기 위한 본 연구의 목적을 밝힘.}
\item{논문의 구성을 서술 (optional).}
\end{itemize}
서론은 과거부터 현재까지 해당 분야의 연구 진행을 기술하기 때문에 선행 연구 논문들을 레퍼런스로 도입하는 경우가 빈번하게 나타난다. \LaTeX에서 참고문헌을 표기하는 방법을 알아보자. 먼저 이 문서의 후반부에 위치한 레퍼런스 부분을 찾아간다. 이 문서를 컴파일했을 때 생성된 PDF 파일에는 {\bf References}라고 나와 있지만 여기서는 {\textbackslash}begin\{thebibliography\}\{99\}로 시작에서 {\textbackslash}end\{thebibliography\}로 종료되는 그 사이에 참고문헌을 작성하면 된다. 여기서 숫자 99는 참고문헌이 100개 넘는 논문을 작성하는 것이 아니라면 그대로 놔둔다. 참고문헌 작성 예시는 다음과 같다.
\begin{lstlisting}
\bibitem{Mok06}C. Mok, C.-M. Ryu, P. H. Yoon, and A. T. Y. Lui, ``Global twofluid stability of bifurcated current sheets'', J. Goephys. Res., {\bf 111}, A03203
(2006).
\end{lstlisting}
{\textbackslash}bibitem 다음의 \{ \} 안에는 자신이 그 논문을 기억하기 쉬운 규칙을 정하여 작성하면 된다. 보통 논문 주저자의 last name과 논문 출판 년도를 사용하여 표기한다. 그리고 저자들, 논문 제목, 저널 이름, 권, 호, 페이지, 출판년도 순으로 입력한다. 저자는 3인 이하일 경우에는 모두 적도록 하고 4인 이상일 때는 주저자만 작성하고 그 외는 et al.이란 표기로 대체한다. 논문 제목은 큰 따옴표로 묶어준다. \LaTeX에서 시작하는 따옴표는 키보드에서 숫자 1 왼쪽 버튼, 마치는 따옴표는 키보드 엔터키 왼쪽 버튼을 사용함에 유의하라. 저널 이름은 경우에 따라 약어를 사용할 수 있다. 약어를 사용할 때는 정식으로 정해진 약어임을 확인한 후 사용한다. Volume(권)은 굵은 글자 처리한다. 위의 예시처럼 {\textbackslash}bf를 사용하면 된다. Number(호)는 경우에 따라 없는 저널도 있다. 위의 예시에는 `호'를 넣지 않았다. 만약 `호'를 넣고자 할 때는 둥근 괄호로 묶어준다. 마지막으로 페이지와 출판년도를 작성한다. 출판년도는 둥근 괄호로 묶어준다.

이제 서론에서 해당 논문을 인용할 준비 작업은 끝났다. 서론에서 필요한 부분에 이 논문을 인용 표기할 경우 {\textbackslash}cite 라고 입력한 후 \{ \} 안에 해당 논문을 표시하면 된다. 표시하는 방법은 바로 레퍼런스에서 {\textbackslash}bibitem 이후 \{ \} 안에 적었던 것을 넣어주면 된다. 논문 인용 표시가 문장 마지막에 등장할 때는 마침표의 위치는 인용 표시 다음이다. 아래 문장은 논문 인용 표시의 예로 C. Mok의 2010년 논문에서 인용하였다 \cite{Mok10}.
\begin{lstlisting}
Various plasma instabilities have been proposed as playing important roles during the substorm onset process. These include the tearing \cite{Schindler74, Sitnov97, Zelenyi08}, ballooning \cite{Cheng98, Bhattacharjee98, Dobias04, Zhu03, Saito08, Friedrich01}, lower hybrid drift \cite{Shinohara98, Yoon02, Mok06}, Kelvin--Helmholtz \cite{Rostoker84, Dovias06}, and the ion Weibel \cite{Yoon93, Sadovskii01} instabilities.
\end{lstlisting}
위와 같이 입력한 후 컴파일하면 pdf 파일에는 다음과 같이 나타날 것이다.
\begin{quote}
Various plasma instabilities have been proposed as playing important roles during the substorm onset process. These include the tearing \cite{Schindler74, Sitnov97, Zelenyi08}, ballooning \cite{Cheng98, Bhattacharjee98, Dobias04, Zhu03, Saito08, Friedrich01}, lower hybrid drift \cite{Shinohara98, Yoon02, Mok06}, Kelvin--Helmholtz \cite{Rostoker84, Dovias06}, and the ion Weibel \cite{Yoon93, Sadovskii01} instabilities.
\end{quote}
이 때 참고문헌이 번호 순서대로 나오도록 한다. 또한 세 개 이상의 문헌이 연속된 번호로 이어진 경우 자동으로 첫 번호와 마지막 번호가 hyphen으로 연결된 형태로 등장함을 확인할 수 있다.


참고문헌은 다음의 조건들을 만족해야 한다.
\begin{itemize}
\item{저자가 명시되어야 한다.}
\item{검증이 된 내용이어야 한다.}
\item{이미 출판되어 수정이 불가능해야 한다.}
\end{itemize}
전문 논문 저널에 수록된 논문들은 위 조건들을 만족하므로 되도록 논문을 참고문헌으로 삼도록 한다. 웹사이트는 위 조건들을 만족하지 못하므로 참고문헌으로 부적절하다. 또한 누구라도 책을 출판할 수 있으므로 전문 서적을 참고문헌으로 사용하는 경우에는 널리 받아들여지고 인정받는 서적만 사용해야 한다. 사실 전공 서적의 저자는 여러 연구 논문들을 참고로 하여 책을 집필하기 때문에 전공서적에도 참고 문헌(논문)이 명시되어 있다. 이 경우 전공 서적 대신에 책에서 지시하는 논문을 참고문헌으로 삼도록 한다.

 % 서론
	\section{Theoretical Background}

\subsection{줄 바꿈}

\LaTeX 에서 줄 바꿈을 하는 방법은 한컴이나 워드와는 다르다.
줄 바꿈을 하기 위해서는 Enter 키를 2번 연속으로 누르거나, 
Back slash 키를 두 번 연속 누르는 방법(\textbackslash \textbackslash)을 사용한다. 
또한 여러 줄을 띄우고 싶다면, \textbackslash vskip 명령어를 사용한다. 
예를 들어, \textbackslash vskip 2pc을 사용 시
\vskip 2pc
위와 같이 한 번에 두 줄을 띄운다. (pc 대신 cm나 mm 등도 사용 가능하다)


그런데 \textbackslash \textbackslash 를 제외한 나머지 방법 사용 시 첫 단어에 들여쓰기가 되는 것을 볼 수 있다. 들여쓰기를 하지 않으려면 문단의 첫부분에 \textbackslash noindent 명령어를 사용한다.

\subsection{그림 삽입}

\LaTeX 에서 그림을 넣는 명령어는 아래와 같다.
\begin{lstlisting}
	\begin{figure}[t]
		\begin{center}
			\includegraphics[width=.3\textwidth]{figure.png}
			\caption{Figure Caption}
			\label{figlabel}
		\end{center}
	\end{figure}
\end{lstlisting}
그림을 넣기 위해서는 images 폴더에 그림이 있어야 한다.
위 명령어에서 width는 그림의 너비, caption은 그림 밑에 표시될 설명, label은 그림을 언급할 때 사용하는 그림의 코드이다. 그림의 언급은 \textbackslash ref\{figlabel\}과 같은 방법으로 한다.

\noindent center 옆의 [t]는 그림을 넣을 위치이다. h, t, b, p 총 네 가지의 설정이 가능하며 각각 here(여기), top(페이지 맨 위), bottom(페이지 맨 아래), page(새로운 페이지)를 의미한다.
일반적으로 htbp를 사용하지만 논문과 같이 페이지 상단에 그림을 위치시켜야 하는 경우 [t]를 사용한다.

\newpage

\begin{figure}[t]
	\begin{center}
		\includegraphics[width=.3\textwidth]{figure.png}
		\caption{Figure Caption}
		\label{figlabel}
	\end{center}
\end{figure}

\noindent 예시 코드 실행 시 위와 같이 페이지 상단에 그림 \ref{figlabel}이 나타난다.

\subsection{표 삽입}

\LaTeX 에서 직접 표를 만들 수 있으나 명령어가 길고 복잡하며, 표의 모습을 보면서 작업할 수 없기에 본인이 원하는 모양대로 만드는 것이 불편하다는 단점이 있다.
그러므로 표의 경우 다른 조판 프로그램에서 제작한 후 이미지 형태로 저장해 \LaTeX 문서에 넣거나, 인터넷 사이트인 Tables Generator 등을 사용하는 것이 더 효율적이다. 다음은 Tables Generator를 사용해 만든 표이다.
\begin{lstlisting}
\begin{table}[htbp]
	\centering
	\caption{렌즈의 종류}
	\label{lens}
	\resizebox{.5\textwidth}{!}{%
		\begin{tabular}{|l|l|l|}
			\hline
			\multicolumn{1}{|c|}{} & 광각렌즈 & 망원렌즈 \\ \hline
			초점거리                   & 짧다   & 길다   \\ \hline
			화각                     & 크다   & 작다   \\ \hline
		\end{tabular}%
	}
\end{table}
\end{lstlisting}

\begin{table}[htbp]
	\centering
	\caption{렌즈의 종류}
	\label{lens}
	\resizebox{.5\textwidth}{!}{%
		\begin{tabular}{|l|l|l|}
			\hline
			\multicolumn{1}{|c|}{} & 광각렌즈 & 망원렌즈 \\ \hline
			초점거리                   & 짧다   & 길다   \\ \hline
			화각                     & 크다   & 작다   \\ \hline
		\end{tabular}%
	}
\end{table}

\subsection{수식 삽입}

문장 내부에서 수식을 사용하기 위해서는 \$ 수식 \$이나 \textbackslash( 수식 \textbackslash)와 같은 방법을 사용한다.
예를 들어 '소비 전력 \$P=I\^{}2 R\$이다.'라는 코드를 작성 후 컴파일하면
\begin{quote}
'소비 전력 $P=I^2 R$이다.'
\end{quote}
\noindent 와 같이 수식이 나타난다. \\
이때 수식과 문장의 높이를 맞추기 위해 분수나 시그마 등의 수학 기호의 크기가 조절된다. 수식의 가독성을 키우기 위해서는 수식 앞에 \textbackslash displaystyle을 붙이면 된다.
\begin{itemize}
	\item displaystyle을 사용하지 않은 경우 : $\sum_{i=1}^{n} = \frac{n(n+1)}{2}$
	\item displaystyle을 사용한 경우 : \(\displaystyle \sum_{i=1}^{n} = \frac{n(n+1)}{2}\)
\end{itemize}

\vskip 2pc

\noindent 수식이 문장 안에 있는 것이 아니라 한 줄을 통째로 차지하는 경우 \textbackslash[ 수식 \textbackslash]이나 명령어 \textbackslash equation을 사용한다. 차이점이라면 전자는 수식에 번호가 붙지 않고, 후자는 붙는다는 것이다. 각각을 사용한 명령어와 실행 모습은 다음과 같다.

\begin{lstlisting}
\[ K + U = \frac{1}{2} mv^2 - \frac{GMm}{r^2}\]
\end{lstlisting}
\[ K + U = \frac{1}{2} mv^2 - \frac{GMm}{r^2}\]

\vskip 2pc

\noindent \begin{lstlisting}
\begin{equation}
S = k_B \ln \Omega \label{equ1}
\end{equation}
\end{lstlisting}
\begin{equation}
	S = k_B \ln \Omega \label{equ1}
\end{equation}

\noindent \textbackslash equation을 사용해서 작성한 식을 언급하기 위해서는 \textbackslash ref\{equ1\}과 같이 ref 명령어를 사용한다. \\
예시 : 통계역학적 엔트로피는 식 \ref{equ1}(\textbackslash ref\{equ1\})과 같이 표현된다.

\vskip 2pc

\noindent 수식이 길어서 한 줄에 표현할 수 없는 경우에는 align 명령어를 사용한다.
\begin{lstlisting}
\begin{align}
F &= \frac{dp}{dt} \notag\\
&= \frac{d}{dt}(mv) \notag \\
&= \frac{dm}{dt}a + m\frac{dv}{dt} \notag \\
&= ma \textrm{ (if mass is constant)} \label{equ2}
\end{align}
\end{lstlisting}
위 명령어를 입력한 후 컴파일하면

\begin{align}
F &= \frac{dp}{dt} \notag\\
&= \frac{d}{dt}(mv) \notag \\
&= \frac{dm}{dt}a + m\frac{dv}{dt} \notag \\
&= ma \textrm{ (if mass is constant)} \label{equ2}
\end{align}

\noindent 위와 같이 여러 줄로 이루어진 수식 \ref{equ2}가 나온다. \\
등호 앞에 있는 \& 기호는 \&를 기준으로 줄을 맞추라는 의미이며 \\
중간 과정의 수식에는 번호를 붙이지 않게 하기 위해 마지막 줄을 제외한 각 줄의 끝에 \textbackslash notag를 붙인다.\\ % 이론적 배경
	\section{연구 과정}
%\section{Method of the study}

\subsection{제목}

내용 % 연구 과정
	\section{결과 및 토의}
%\section{Result and discussion}

\subsection{제목}

내용 % 결과 및 토의
	\section{결론}
%\section{Conclusion}

\subsection{제목}

내용 % 결론
	% \section{추후 연구}
%\section{Further research}

\subsection{제목}

내용 % 추후 연구
	\section{부록}
%\section{Appendix} % 부록
	
	\bibliography{bibfile} % 참고문헌
	% BibTeX 코드 쉽게 얻어오는 방법 %
	% Google Scholar 에서 검색한 결과에서 `인용'을 클릭한다.
	% BibTeX 코드를 얻고자 한다면, 하단의 `BibTeX' 을 클릭.
	% 코드가 나온다. Ctrl+A, Ctrl+C로 복사, bibfile에 붙여넣기.


\end{document}
